\documentclass[11pt,a4paper]{article}
\usepackage[margin=1in]{geometry}
\usepackage{amsmath,booktabs,hyperref,float}

\title{Iteration 010: Deep Temporal Convolution}
\author{Autoresearch Loop}
\date{April 7, 2026}

\begin{document}
\maketitle

\begin{abstract}
A deeper convolutional model with depthwise-separable convolutions and residual
connections achieves $r=0.372$, matching but not exceeding the simpler FIR filter
($r=0.373$). The added capacity does not improve over the linear spatio-temporal
baseline, confirming that the 1--9\,Hz mapping is approximately linear.
\end{abstract}

\section{Hypothesis}
Stacking temporal convolutions with nonlinear activations could capture nonlinear
temporal dynamics that the linear FIR misses. Depthwise-separable convolutions
reduce parameter count while maintaining a larger receptive field.

\section{Method}
Architecture: depthwise 1D conv (per-channel temporal filtering) followed by
pointwise 1D conv (cross-channel mixing), with residual connections and ReLU
activations. Trained 100 epochs with AdamW, CF center-tap initialization.

\section{Results}
\begin{table}[H]
\centering
\begin{tabular}{lrrr}
\toprule
Model & Mean $r$ & Std $r$ & SNR (dB) \\
\midrule
FIR baseline (iter009) & 0.373 & 0.074 & 0.61 \\
Deep temporal conv & 0.372 & 0.076 & 0.62 \\
\bottomrule
\end{tabular}
\end{table}

\section{Analysis}
The nonlinear model offers no improvement. In the 1--9\,Hz band at 20\,Hz,
the scalp-to-inear relationship is well-captured by a linear spatio-temporal
filter. Adding depth and nonlinearity introduces more parameters to overfit
without learning useful nonlinear features. Future iterations should focus
on regularization and robustness rather than model capacity.

\section{Next}
Iteration 011: channel dropout augmentation to improve cross-subject robustness.

\end{document}
